\documentclass[10pt, t, aspectratio = 1610]{beamer}
\usepackage{./template/beamertheme}
\usepackage{amsmath, bm, mathtools}
  \newcommand{\e}{\mathrm{e}}
  \newcommand{\dif}{\mathrm{d}}
  \DeclareMathOperator{\relu}{relu}
\usefonttheme{professionalfonts}

\makeatletter
\colorlet{barcolor}{beamer@barcolor}
\colorlet{headercolor}{beamer@headercolor}
\makeatother
\usepackage[UTF8, linespread = 1.5]{ctex}
\usepackage{multicol}
\setlength{\columnsep}{1cm}
\usepackage{lipsum}
\usepackage{parskip}
\usepackage{contour}
  \contourlength{1pt}
  \contournumber{50}
  \newcommand{\shadowtext}[1]{\contour{white}{#1}}
\usepackage[many, minted]{tcolorbox}
\tcbuselibrary{breakable}
\makeatletter
\renewenvironment{quote}{%
    \begin{tcolorbox}%
    [
        top = 3pt,
        bottom = 3pt,
        left = 3pt,
        right = 3pt,
        colback = beamer@barcolor!10!white,
        colframe = beamer@headercolor,
        colupper = beamer@normaltextcolor,
    ]%
}%
{\end{tcolorbox}}

\def\inlinebox@true{true}
\define@key{inlinebox}{frame rule color} {\def\inlinebox@framerulecolor{#1}}
\define@key{inlinebox}{frame back color} {\def\inlinebox@framebackcolor{#1}}
\define@key{inlinebox}{frame text color} {\def\inlinebox@frametextcolor{#1}}
\define@key{inlinebox}{frame rule width} {\def\inlinebox@framerulewidth{#1}}
\define@key{inlinebox}{banner width}     {\def\inlinebox@bannerwidth{#1}}
\define@key{inlinebox}{show banner}[true]{\def\inlinebox@showbanner{#1}}
\define@key{inlinebox}{banner text color}{\def\inlinebox@bannertextcolor{#1}}
\define@key{inlinebox}{banner back color}{\def\inlinebox@bannerbackcolor{#1}}
\define@key{inlinebox}{banner text}      {\def\inlinebox@bannertext{#1}}
\NewDocumentCommand{\inlinebox}{O{} m}{%
  \setkeys{inlinebox}{%
    frame rule color  = black,
    frame back color  = white,
    frame text color  = black,
    frame rule width  = 0.4pt,
    banner width      = 8pt,
    show banner       = false,
    banner text color = white,
    banner back color = black,
    banner text       = BAN,
    #1
  }%
  \tcbox[%
    enhanced,
    tcbox raise base,
    nobeforeafter,
    boxrule           = \inlinebox@framerulewidth,
    top               = -1pt,
    bottom            = -1pt,
    right             = -1pt,
    arc               = 1pt,
    left              = \ifx\inlinebox@showbanner\inlinebox@true\inlinebox@bannerwidth-2pt\else-1pt\fi,
    colframe          = \inlinebox@framerulecolor,
    coltext           = \inlinebox@frametextcolor,
    colback           = \inlinebox@framebackcolor,
    before upper      = {\vphantom{蛤dg}},
    overlay           = {%
      \begin{tcbclipinterior}
        \ifx\inlinebox@showbanner\inlinebox@true
          \fill[\inlinebox@bannerbackcolor] (frame.south west) rectangle node[text = \inlinebox@bannertextcolor, scale = 0.4, font = \sffamily\bfseries, rotate = 90] {\inlinebox@bannertext} ([xshift = \inlinebox@bannerwidth]frame.north west);
        \fi
      \end{tcbclipinterior}%
    }%
  ]{#2}%
}

\newtcblisting{pythoncode}[1][]{
    listing engine = minted,
    minted style = colorful,
    minted language = python,
    minted options = {%
        fontsize = \fontsize{9}{9}\selectfont,
        linenos,
        mathescape = true,
        numbersep = 3mm,
        autogobble,
        breaklines,
        breakafter = d,
        #1},
    colback = barcolor!5!white,
    colframe = headercolor,
    listing only,
    left = 5mm,
    top = 0pt,
    bottom = 0pt,
    enhanced,
    overlay = {%
        \begin{tcbclipinterior}\fill[headercolor!20!white] (frame.south west)
            rectangle ([xshift=5mm]frame.north west);\end{tcbclipinterior}}%
}

\newtcblisting{shellcode}[1][]{
    listing engine = minted,
    minted style = colorful,
    minted language = text,
    minted options = {%
        fontsize = \fontsize{9}{9}\selectfont,
        linenos,
        escapeinside=||,
        mathescape = true,
        numbersep = 3mm,
        autogobble,
        breaklines,
        breakafter = d,
        #1},
    colback = barcolor!5!white,
    colframe = headercolor,
    listing only,
    left = 5mm,
    top = 0pt,
    bottom = 0pt,
    enhanced,
    overlay = {%
        \begin{tcbclipinterior}\fill[headercolor!20!white] (frame.south west)
            rectangle ([xshift=5mm]frame.north west);\end{tcbclipinterior}}%
}
% \newtcblisting{pythoncode}{
%     colback = beamer@barcolor!5!white,
%     colframe = beamer@headercolor,
%     left = 5mm,
%     top = 0pt,
%     bottom = 0pt,
%     enhanced,
% }

\newcommand{\inlinepython}[1]{\inlinebox[frame rule color = black!30!white]{\mintinline{python}{#1}}}
\newcommand{\inlinetext}[1]{\inlinebox[frame rule color = black!30!white]{\mintinline{text}{#1}}}
\newcommand{\inlineframe}[1]{\inlinebox[frame rule color = black!30!white]{#1}}
\makeatother



\usepackage{enumitem}
  \setitemize[1]{itemsep = 0pt, partopsep = 0pt, parsep = 0pt, topsep = 0pt, leftmargin = 1em, label = $\bullet$}
  \setitemize[2]{itemsep = 0pt, partopsep = 0pt, parsep = 0pt, topsep = 0pt, leftmargin = 1em, label = $-$}
  \usepackage{tabu}
  \tabuDisableCommands{%
    \renewcommand\mintinline[2]{\texttt{\detokenize{#2}}}%
  }
  \tabulinesep = 3pt
\makeatletter
  \setlength{\parskip}{\medskipamount}

  \newcommand{\@minipagerestore}{\setlength{\parskip}{\medskipamount}}
  \makeatother
\title{机器学习入门}
\subtitle{以及如何使用 TensorFlow}
\date{\today}
\author{%
    张琦%
}

\institute%
{%
    2012 实验室 $\cdot$ 测试工具部\\%
    华为技术有限公司%
}

\pgfdeclareimage[height = 1cm]{titlepagelogo}{./template/icon_in_cover.pdf}
\titlegraphic{%
  \pgfuseimage{titlepagelogo}%
}

\newcommand{\insertcover}{{
    \aauwavesbg%
    \begin{frame}[plain, noframenumbering]
         \titlepage
    \end{frame}
}}

\newcommand{\insertthanks}{{
    \aauwavesbg
    \begin{frame}[plain, noframenumbering]
        \finalpage{\huge 谢谢, 欢迎提问\vspace{-10pt}}
    \end{frame}
}}

\usepackage{pgfplots}
\usetikzlibrary{arrows.meta}
\tikzset
{
  > = {Stealth[length = 5pt]},
}

\usepgfplotslibrary{groupplots}
\pgfplotsset{
    compat = newest,
    discard if not/.style 2 args={
        x filter/.code={
            \edef\tempa{\thisrow{#1}}
            \edef\tempb{#2}
            \ifx\tempa\tempb
            \else
                \def\pgfmathresult{inf}
            \fi
        }
    },
    every tick label/.append style = {font=\scriptsize},
    label style = {inner sep = 0pt, font=\footnotesize},
    title style = {inner sep = 0pt, font=\footnotesize},
    legend style={nodes={scale=0.6, transform shape}},
    layers/my layer set/.define layer set={
        background,
        main,
        foreground
    }{
    },
    set layers=my layer set,
}
\setbeamertemplate{section in toc}[sections numbered]
\setbeamertemplate{subsection in toc}[subsections numbered]

\usepackage{varwidth}
\usepackage{etoolbox}
\makeatletter
%\patchcmd{\beamer@sectionintoc}{\vskip1.5em}{\vskip0.5em}{}{}
\patchcmd{\beamer@sectionintoc}{%
  \hbox{\vbox{%
    \def\beamer@breakhere{\\}%
    \beamer@tocact{\ifnum\c@section=#1\beamer@toc@cs\else\beamer@toc@os\fi}{section in toc}}}%
}{%
  \hbox{%
    \def\beamer@breakhere{}%
    \beamer@tocact{\ifnum\c@section=#1\beamer@toc@cs\else\beamer@toc@os\fi}{section in toc}}%
}{}{}
\makeatother

\newcommand{\tensorflow}{\tikz[baseline]{\node[anchor = base, inner sep = 0pt] {\raisebox{-1pt}{\includegraphics[height = 1.1em]{./template/icon_in_cover.pdf}}};}}
\begin{document}
\insertcover

% \begin{frame}{目录}
%     \tableofcontents[hideallsubsections]
% \end{frame}

%!TEX root = ../main.tex

\section{开胃菜}
\begin{frame}{\secname}
\vfill\hspace{20pt}
\begin{minipage}[m]{0.35\textwidth}
  \linespread{1}\large
  \tableofcontents[sectionstyle = hide/hide, subsectionstyle = show/show/hide]
\end{minipage}%
\hfill%
\begin{minipage}[m]{0.5\textwidth}
  \includegraphics[height = 7cm]{./icons/appetizer.pdf}
\end{minipage}%
\vfill
\end{frame}

\subsection{简单的分类问题}
\begin{frame}{\insertsection}{\insertsubsection}
\begin{minipage}[m]{0.48\textwidth}
在二维空间中, 有 $500$ 个点, 这 $500$ 个点被分成两类, 一类被标记为绿色, 一类被标记为蓝色, 如右图所示.

现在的问题是, 计算机如何根据已有的数据, 找到一个分类的方法, 可以在给定 $(x_1, x_2)$ 时, 会判断点 $(x_1, x_2)$ 的颜色.

计算机解决二分类的方法有很多, 在这里, 介绍一种基于神经网络的方法, 来解决这个问题, 来帮助大家理解机器学习的过程.
\end{minipage}%
\hfill%
\begin{minipage}[m]{0.48\textwidth}
\begin{tikzpicture}
  \begin{axis}[%
    xlabel = $x_1$,
    ylabel = $x_2$,
    height = 7.8cm,
    width = 7cm,
    title = $500$ 个点的分布图
  ]
    \addplot[only marks, discard if not={type}{1}, blue] table[col sep=comma, x = x1, y = x2]{./data/classification/training_data.dat};
    \addplot[only marks, discard if not={type}{0}, green] table[col sep=comma, x = x1, y = x2]{./data/classification/training_data.dat};
  \end{axis}
\end{tikzpicture}
\end{minipage}
\end{frame}

\subsection{建立神经网络}
\begin{frame}{\insertsection}{\insertsubsection}
我们不妨设绿色为 $1$, 蓝色为 $0$, 那么, 问题就变成了找到一个坐标 $(x_1, x_2)$ 到颜色 $y$ 的一个映射关系 $y = h(x_1, x_2)$.\vspace{10pt}

\begin{minipage}[m]{0.55\textwidth}
为此, 我们建立一个简单的神经网络, 如图所示.
\begin{center}
\scalebox{0.8}{%
\begin{tikzpicture}[%
  node/.style = {circle, draw, inner sep = 0pt, minimum size = 0.8cm},
  x = 1.4cm]
  \node[node] (x1) at (0, 2) {$x_1$};
  \node[node] (x2) at (0, 1) {$x_2$};
  \node[node] (b)  at (0, 0) {$1$};
  \node[node] (sum) at (2, 1) {$\sum$};
  \node[node] (step) at (4, 1) {$H(x)$};

  \draw[->] (x1) -- (sum) node[midway, fill = white, inner sep = 1pt, sloped] {$W_1$};
  \draw[->] (x2) -- (sum) node[midway, fill = white, inner sep = 1pt, sloped] {$W_2$};
  \draw[->] (b) -- (sum);
  \draw[->] (sum) -- (step);
  \draw[->] (step) -- ++(0:1.5cm) node[anchor = west] {$y$};
\end{tikzpicture}}
\end{center}
因此, 我们可以得到神经网络的表达式:\vspace{-10pt}
\begin{align*}
  y &= H(W_1x_1 + W_2x_2 + 1)\text{,}\\
    &= H\big((W_1, W_2)\cdot(x_1, x_2)^\top + 1\big) = H(\bm{W}\bm{x} + 1)\text{.}
\end{align*}
\end{minipage}%
\hfill%
\begin{minipage}[m]{0.43\textwidth}
\begin{tikzpicture}
  \begin{axis}[width = \textwidth,
               height = 5.7cm,
               xlabel = $x$,
               ylabel = $H(x)$,
               title = 单位阶跃函数图像]
    \addplot[mark = none, samples = 100, blue] {x > 0};
  \end{axis}
\end{tikzpicture}
\end{minipage}
\end{frame}

\subsection{确定目标函数}
\begin{frame}{\insertsection}{\insertsubsection}
所谓目标函数, 就是当目标函数取得最值的时候, 神经网络最符合我们的预期. 换句话就是, 目标函数就是衡量神经网络好坏指标.

我们这里采用所有样本都分类成功的概率作为目标函数:
\[
p(X) = \prod_{i = 1}^{500} p(X_i)\text{,}
\]

对于第 $i$ 个坐标 $\bm{x}_i = (x_{i, 1}, x_{i, 2})$,其正确值是 $y_i$,那么当曲线 $h(x_1, x_2)$ 给定的情况下,其作出正确分类的概率为%
%
\begin{align*}
  p(X_i) &= \left\{
  \begin{array}{ll}
    1, & \text{$H(\bm{W}\bm{x}_i + 1) = 1, y_i = 1$ 或者 $H(\bm{W}\bm{x}_i + 1) = 0, y_i = 0$,}\\
    0, & \text{$H(\bm{W}\bm{x}_i + 1) = 0, y_i = 1$ 或者 $H(\bm{W}\bm{x}_i + 1) = 1, y_i = 0$.}
  \end{array}
  \right.
\end{align*}
\end{frame}

\begin{frame}{\insertsection}{\insertsubsection}
接下来, 需要对这个式子进行化简.
\[
  p(X_i) = \left\{
  \begin{array}{ll}
    1, & \text{$H(\bm{W}\bm{x}_i + 1) = 1, y_i = 1$ 或者 $H(\bm{W}\bm{x}_i + 1) = 0, y_i = 0$,}\\
    0, & \text{$H(\bm{W}\bm{x}_i + 1) = 0, y_i = 1$ 或者 $H(\bm{W}\bm{x}_i + 1) = 1, y_i = 0$.}
  \end{array}
  \right.
\]

我们可以注意到, 当 $y_i = 1$ 时, $p(X_i)$ 的值与 $H(\bm{W}\bm{x}_i + 1)$ 一致; 当 $y_i = 0$ 时, $p(X_i)$ 的值与 $H(\bm{W}\bm{x}_i + 1)$ 和为 $1$, 于是有:
\[
  p(X_i) = \left\{
  \begin{array}{ll}
    H(\bm{W}\bm{x}_i + 1), & y_i = 1\text{,}\\
    1 - H(\bm{W}\bm{x}_i + 1), & y_i = 0\text{.}
  \end{array}
  \right.
\]

注意到 $\forall x \in \mathbb{R}$ 都有 $x^0 = 1$, 还可以将上式继续化简为:
\[
  p(X_i) = H(\bm{W}\bm{x}_i + 1)^{y_i}\cdot\big(1 - H(\bm{W}\bm{x}_i + 1)\big)^{1 - y_i}\text{.}
\]
\end{frame}

\begin{frame}{\insertsection}{\insertsubsection}
因此, 我们的目标函数就确定下来了:
\[
  p(X) = \prod_{i = 1}^{500} H(\bm{W}\bm{x}_i + 1)^{y_i}\cdot\big(1 - H(\bm{W}\bm{x}_i + 1)\big)^{1 - y_i}\text{.}
\]

为了方便计算, 对上式左右两边都取以 $\e$ 为底的对数:
\[
  \ln p(X) = \sum_{i = 1}^{500} y_i\ln H(\bm{W}\bm{x}_i + 1) + (1 - y_i)\ln\big(1 - H(\bm{W}\bm{x}_i + 1)\big)\text{.}
\]

但是, 上述目标函数还存在问题:%
%
\begin{itemize}
\item 导致函数 $\ln p(X)$ 不可微, 不能采用梯度下降法来计算函数的最值(其实是极值);
\item 只要有一个预测错误, 函数 $\ln p(X) = -\infty$, 目标函数并不能衡量模型的精度.
\end{itemize}
\end{frame}

\begin{frame}{\insertsection}{\insertsubsection}
解决方法就是用 Sigmoid 函数 $g(\,\cdot\,)$ 代替上式中的单位阶跃函数 $H(\,\cdot\,)$, 原因有:
\begin{itemize}
\item Sigmoid 函数与单位阶跃函数形状非常相似, 并且值域为 $(0, 1)$, 不会导致目标函数出现 $-\infty$ 的情况, 其函数图像如右图所示.
\end{itemize}\vspace{-10pt}
\begin{minipage}[t]{0.48\textwidth}
\begin{itemize}
\item Sigmoid 处处可微, Sigmoid 的表达式为:%
\[
  g(x) = \frac{1}{1 + \exp(-x)}\text{.}
\]
\item Sigmoid 函数有非常好的导函数性质, 有:%
\[
  \frac{\dif g(x)}{\dif x} = g(x)\cdot\big(1 - g(x)\big)\text{.}
\]
\end{itemize}
\end{minipage}%
\hfill%
\begin{minipage}[t]{0.48\textwidth}
\begin{figure}
  \centering\vspace{-20pt}
  \begin{tikzpicture}
    \begin{axis}[height = 4.5cm,
                 width = \textwidth,
                 xlabel = $x$,
                 title = Sigmoid 函数与单位阶跃函数图像,
                 legend pos = north west,]
      \addplot[mark = none, blue, domain = -30:30, samples = 100] {(x > 0)};
      \addplot[mark = none, green, domain = -30:30, samples = 100] {(x <= 0) * exp(x)/(1 + exp(x)) + (x > 0) / (1 + exp(-x))};
      \legend{单位阶跃函数, Sigmoid 函数};
    \end{axis}
  \end{tikzpicture}
\end{figure}
\end{minipage}
\begin{flalign*}
& \text{目标函数:} &&  \ell(\bm{W}) = -\sum_{i = 1}^{500} y_i\ln g(\bm{W}\bm{x}_i + 1) + (1 - y_i)\ln\big(1 - g(\bm{W}\bm{x}_i + 1)\big)\text{.} &&
\end{flalign*}
\end{frame}

\subsection{参数训练}
\begin{frame}{\insertsection}{\insertsubsection}
参数训练的本质是找到一个参数 $\bm{W}$ 使目标函数 $\ell(\bm{W})$ 取得最小值, 常用的是梯度下降法.

\begin{quote}
把 $\bm{W} = (W_1, W_2)$ 看成是平面坐标, $\ell(\bm{W})$ 看作是坐标 $(W_1, W_2)$ 的海拔高度, 梯度下降法的策略就是每次向最陡的方向走一步, 一直走到海拔不变为止, 此时就是海拔``最低点'', 即 $\ell(\bm{W})$ 的极小值点.
\end{quote}\vspace{-15pt}

\tikzstyle{every picture}+=[remember picture]
\tikzstyle{na} = [baseline=-.5ex]
\begin{flalign*}
& \text{梯度下降的迭代公式为: } &&
\tikz[baseline]{\node[draw = blue, anchor = base, rounded corners] (w1) {$\bm{W}_n$};} \gets
\tikz[baseline]{\node[draw = red, anchor = base, rounded corners] (w2) {$\bm{W}_{n-1}$};} -
\tikz[baseline]{\node[draw = black, anchor = base, rounded corners] (r) {$\bm{r}$};}\cdot
\tikz[baseline]{\node[draw = green, anchor = base, rounded corners] (p) {$\dfrac{\partial \ell(\bm{W})}{\partial \bm{W}}\Big|_{\bm{W} = \bm{W}_{n-1}}$};}
\text{,}
&& \text{其中:}
\end{flalign*}

\begin{itemize}
\item 下一步的位置;\tikz[]{\node[anchor = base] (w1') {\ };}
\item 当前站的位置;\tikz[]{\node[anchor = base] (w2') {\ };}
\item 步长, 也被称为学习率;\tikz[]{\node[anchor = base] (r') {\ };}
\item 最陡的方向, 也就是梯度.\tikz[]{\node[anchor = base] (p') {\ };}
\end{itemize}

\begin{tikzpicture}[overlay]
\draw[->, rounded corners, blue] (w1') -| ([yshift = -2pt]w1.south);
\draw[->, rounded corners, red] (w2') -| ([yshift = -2pt]w2.south);
\draw[->, rounded corners, black] (r') -| ([yshift = -2pt]r.south);
\draw[->, rounded corners, green] (p') -| ([yshift = -2pt]p.south);
\end{tikzpicture}
\end{frame}

\begin{frame}{\insertsection}{\insertsubsection}
我们计算一下 $\frac{\partial \ell(\bm{W})}{\partial \bm{W}}$, 设 $z = \bm{W}\bm{x} + 1 = W_1x_{i,1} + W_2x_{i,2} + 1$, 于是:
\begin{align*}
\dfrac{\partial \ell(\bm{W})}{\partial z} &= -\dfrac{\partial \displaystyle\sum_{i = 1}^{500} \Big(y_i\ln g(z) + (1 - y_i)\ln\big(1 - g(z)\big)\Big)}{\partial z}\text{,}\\
  &= -\sum_{i = 1}^{500} y_i \dfrac{\partial \ln g(z)}{\partial z} + (1 - y_i)\dfrac{\partial\ln \big(1 - g(z)\big)}{\partial z} = -\sum_{i = 1}^{500} y_i \dfrac{g'(z)}{g(z)} + (1 - y_i)\dfrac{-g'(z)}{1 - g(z)}\text{,}\\
  &= -\sum_{i = 1}^{500} y_i \dfrac{g(z)(1 - g(z))}{g(z)} + (1 - y_i)\dfrac{-g(z)(1 - g(z))}{1 - g(z)} = -\sum_{i = 1}^{500} y_i - g(z)\text{.}\\
  \dfrac{\partial \ell(\bm{W})}{\partial W_1} &= -\dfrac{\partial \ell(\bm{W})}{\partial z} \cdot \dfrac{\partial z}{\partial W_1} = -\sum_{i = 1}^{500} \big(y_i - g(W_1x_{i,1} + W_2x_{i,2} + 1)\big)\cdot x_{i,1}\text{,}\\
  \dfrac{\partial \ell(\bm{W})}{\partial W_2} &= -\dfrac{\partial \ell(\bm{W})}{\partial z} \cdot \dfrac{\partial z}{\partial W_2} = -\sum_{i = 1}^{500} \big(y_i - g(W_1x_{i,1} + W_2x_{i,2} + 1)\big)\cdot x_{i,2}\text{.}
\end{align*}
\end{frame}

\subsection{代码实现}
\begin{frame}[fragile]{\insertsection}{\insertsubsection}
\begin{columns}
\column{\dimexpr\paperwidth-10pt}\vspace{-20pt}
\begin{pythoncode}[fontsize = \fontsize{8}{8}\selectfont]
from math import exp, log

def sigmoid(x, n = 1): # 定义 Sigmoid 函数 $g = 1/(1 + \exp(x))$
  return 1 / (1 + exp(-n * x)) if x >= 0 else exp(n * x) / (1 + exp(n * x))

def z(x1, x2, w1, w2): # 定义中间变量 $z = W_1x_1 + W_2x_2 + 1$
  return w1 * x1 + w2 * x2 + 1

def l(x1, x2, y, w1, w2): # 定义目标函数 $\ell(\bm{W}) =  -\sum_{i = 1}^{500} y_i\ln g(\bm{W}\bm{x}_i + 1) + (1 - y_i)\ln\big(1 - g(\bm{W}\bm{x}_i + 1)\big)$
  return -sum([y * log(sigmoid(z(x1, x2, w1, w2))) + (1 - y) * log(1 - sigmoid(z(x1, x2, w1, w2))) for x1, x2, y in zip(x1, x2, y)])

def get_delta_w1(x1, x2, y, w1, w2): # 定义 $\frac{\partial \ell(\bm{W})}{\partial W_1} = -\sum_{i = 1}^{500} \big(y_i - g(W_1x_{i,1} + W_2x_{i,2} + 1)\big)\cdot x_{i,1}$
  return -sum([(y - sigmoid(z(x1, x2, w1, w2))) * x1 for x1, x2, y in zip(x1, x2, y)])

def get_delta_w2(x1, x2, y, w1, w2): # 定义 $\frac{\partial \ell(\bm{W})}{\partial W_1} = -\sum_{i = 1}^{500} \big(y_i - g(W_1x_{i,1} + W_2x_{i,2} + 1)\big)\cdot x_{i,2}$
  return -sum([(y - sigmoid(z(x1, x2, w1, w2))) * x2 for x1, x2, y in zip(x1, x2, y)])
\end{pythoncode}
\end{columns}
\end{frame}

\begin{frame}[fragile]{\insertsection}{\insertsubsection}
\begin{columns}
\column{\dimexpr\paperwidth-10pt}\vspace{-20pt}
\begin{pythoncode}[fontsize = \fontsize{8}{8}\selectfont]
with open('training.dat', 'r') as data_file: # 读取训练数据 $x_1$, $x_2$ 和 $y$
  next(data_file)
  data = zip(*[[float(x.strip()) for x in line.split(',')] for line in data_file if not line.strip() == ''])
x1, x2, y = (column for column in data)

w1, w2 = 0, 0 # 初始化变量 $W_1$ 和 $W_2$
r = 0.03 # 设置学习率
precision = 1e-4 # 设置训练的目标精度

while True: # 不断的进行循环迭代, 直至 $W_1$ 和 $W_2$ 的值都不变为止
  delta_w1, delta_w2 = get_delta_w1(x1, x2, y, w1, w2), get_delta_w2(x1, x2, y, w1, w2)
  w1, w2 = w1 - r * delta_w1, w2 - r * delta_w2
  if delta_w1**2 + delta_w2**2 < precision**2:
    break

print('w1 = {w1}, w2 = {w2}'.format(w1 = w1, w2 = w2)) # 打印结果
\end{pythoncode}
\end{columns}
\end{frame}

\begin{frame}{\insertsection}{\insertsubsection}
\begin{minipage}[m]{0.32\textwidth}
在训练的过程中, 我们将每一次迭代的结果记录下来, 并绘制成曲线图, 如右图所示.

\begin{itemize}
\item 每一步的方向都是垂直于等高线, 即梯度的方向, 这就是梯度下降法的由来;
\item 步长越来越小, 这是因为坡度越来越缓;
\item 最终参数 $W_1$ 和 $W_2$ 收敛于 $(1.84, -5.46)$.
\end{itemize}
\end{minipage}%
\hfill%
\begin{minipage}[m]{0.65\textwidth}
\begin{tikzpicture}
  \begin{axis}[%
    axis equal = true,
    width = \textwidth,
    xlabel = $W_1$,
    ylabel = $W_2$,
    title = {参数 $W_1$ 和 $W_2$ 的变化轨迹}]
    \addplot[contour prepared] table[col sep=comma]{./data/classification/contour_n_equal_1.dat};
    \addplot[mark = o, mark size = 1pt] table[col sep=comma, x = w1, y = w2]{./data/classification/parameters_n_equal_1.dat};
  \end{axis}
\end{tikzpicture}
\end{minipage}
\end{frame}

\subsection{结果分析}

\begin{frame}{\insertsection}{\insertsubsection}
\begin{minipage}[m]{0.48\textwidth}
由于 参数 $W_1$ 和 $W_2$ 收敛于 $(1.84, -5.46)$, 我们可以得到一条直线:
\[
  1.84x_1 - 5.46x_2 + 1 = 0\text{.}
\]%
%
如右图所示, 可见, 结果很不理想. 通过分析数据, 我们可以看出, 点 $(1.84, -5.46)$ 处的目标函数值为 $\ell(1.84, -5.46) = 105.11$, 也是非常不理想的.

这是由于用 Sigmoid 函数代替单位阶跃函数时, 给目标函数带来的误差导致的.
\end{minipage}%
\hfill%
\begin{minipage}[m]{0.48\textwidth}
\begin{tikzpicture}
  \begin{axis}[%
    xlabel = $x_1$,
    ylabel = $x_2$,
    height = 7.8cm,
    width = 7cm,
    title = 分类效果
  ]
    \addplot[only marks, discard if not={type}{1}, blue] table[col sep=comma, x = x1, y = x2]{./data/classification/training_data.dat};
    \addplot[only marks, discard if not={type}{0}, green] table[col sep=comma, x = x1, y = x2]{./data/classification/training_data.dat};
    \addplot[mark = none, domain = -1:1, white, line width = 1pt, double = red, double distance = 2pt, on layer = foreground] {1.84/5.46*x + 1.0/5.46} node[midway, sloped, text = black] {\shadowtext{$1.84x_1 - 5.46x_2 + 1 = 0$}};
  \end{axis}
\end{tikzpicture}
\end{minipage}
\end{frame}

\subsection{改进激活函数}
\begin{frame}{\insertsection}{\insertsubsection}
\begin{minipage}[t]{0.41\textwidth}
我们对 Sigmoid 函数表达式做稍加改进, 如下所示.
\[
  g(x, n) = \frac{1}{1 + \exp(n\cdot x)}\text{.}
\]

分别绘制单位阶跃函数, $g(x, 1)$ 和 $g(x, 5)$, 可以看出, 随着 $n$ 的增大, 函数 $g(x, n)$ 越来越接近于单位阶跃函数 $H(x)$.
\end{minipage}%
\hfill%
\begin{minipage}[t]{0.55\textwidth}
\begin{figure}
  \centering\vspace{-20pt}
  \begin{tikzpicture}
    \begin{axis}[height = 6cm,
                 width = \textwidth,
                 xlabel = $x$,
                 title = Sigmoid 函数与单位阶跃函数图像,
                 legend pos = north west,]
      \addplot[mark = none, blue, domain = -10:10, samples = 100] {(x > 0)};
      \addplot[mark = none, green, domain = -10:10, samples = 100] {(x <= 0) * exp(x)/(1 + exp(x)) + (x > 0) / (1 + exp(-x))};
      \addplot[mark = none, red, domain = -10:10, samples = 100] {(x <= 0) * exp(5*x)/(1 + exp(5*x)) + (x > 0) / (1 + exp(-5*x))};

      \legend{单位阶跃函数, $1/\big(1 + \exp(x)\big)$, $1/\big(1 + \exp(5x)\big)$};
    \end{axis}
  \end{tikzpicture}
\end{figure}
\end{minipage}

\begin{flalign*}
& \text{并且可以证明:} &&
 \forall x \in \mathbb{R},\,\,\,\, \lim_{n\rightarrow +\infty} g(x, n) = \lim_{n\rightarrow +\infty}\frac{1}{1 + \exp(n\cdot x)} =  H(x)\text{.} &&
\end{flalign*}
\end{frame}

\begin{frame}{\insertsection}{\insertsubsection}
最终参数 $W_1$ 和 $W_2$ 收敛于 $(0.99, -2.93)$. 分类效果和参数 $W_1$ 和 $W_2$ 的轨迹如下所示.\vspace{10pt}

\begin{tikzpicture}
  \begin{groupplot}[group style={group size = 2 by 1, horizontal sep = 1.5cm}]
  \nextgroupplot[%
    axis equal = true,
    xlabel = $x_1$,
    ylabel = $x_2$,
    width = 0.53\textwidth,
    enlarge x limits = 0.1,
    title = 分类效果
  ]
    \addplot[only marks, discard if not={type}{1}, blue] table[col sep=comma, x = x1, y = x2]{./data/classification/training_data.dat};
    \addplot[only marks, discard if not={type}{0}, green] table[col sep=comma, x = x1, y = x2]{./data/classification/training_data.dat};
    \addplot[mark = none, domain = -1:1, white, line width = 1pt, double = red, double distance = 2pt, on layer = foreground] {0.99/2.93*x + 1.0/2.93} node[midway, sloped, text = black] {\shadowtext{$0.99x_1 - 2.93x_2 + 1 = 0$}};
  \nextgroupplot[%
    axis equal = true,
    width = 0.53\textwidth,
    xlabel = $W_1$,
    ylabel = $W_2$,
    title = {参数 $W_1$ 和 $W_2$ 的变化轨迹}]
    \addplot[contour prepared] table[col sep=comma]{./data/classification/contour_n_equal_10.dat};
    \addplot[mark = o, mark size = 1pt] table[col sep=comma, x = w1, y = w2]{./data/classification/parameters_n_equal_10.dat};
  \end{groupplot}
\end{tikzpicture}%
\end{frame}

\subsection{学习率}
\begin{frame}{\insertsection}{\insertsubsection}
\begin{minipage}[m]{0.31\textwidth}
我们把学习率 $r$ 提高, 重新进行学习, 参数 $W_1$ 和 $W_2$ 的变化轨迹如右图所示.

可以看到, 整个目标函数的最优点仍然是 $(0.99, -2.93)$, 但是, 参数 $W_1$ 和 $W_2$ 无法收敛于最优点.

相对的, 过小的学习率会导致特学习过程非常缓慢. 故选择合适大小的学习率是非常重要的.
\end{minipage}%
\hfill%
\begin{minipage}[m]{0.65\textwidth}
\begin{tikzpicture}
  \begin{axis}[%
    axis equal = true,
    width = \textwidth,
    xlabel = $W_1$,
    ylabel = $W_2$,
    title = {参数 $W_1$ 和 $W_2$ 的变化轨迹}]
    \addplot[contour prepared] table[col sep=comma]{./data/classification/contour_n_equal_10_big_learning_rate.dat};
    \addplot[mark = o, mark size = 1pt] table[col sep=comma, x = w1, y = w2]{./data/classification/parameters_n_equal_10_big_learning_rate.dat};
  \end{axis}
\end{tikzpicture}
\end{minipage}
\end{frame}

\subsection{总结}
\begin{frame}{\insertsection}{\insertsubsection}
让我们总结一下这一小节关键点:%
%
\begin{itemize}
\item 机器学习的本质是重现人认识世界的过程, 具体实现靠的是空间搜索和函数的泛化;
\item 机器学习的一般过程:
\begin{itemize}
\item 数据分析以及预处理;
\item 选择合适的模型;
\item 设定目标函数以及训练参数;
\item 验证以及调整模型.
\end{itemize}
\end{itemize}

但是, 机器学习的过程复杂, 需要计算非常复杂的目标函数, 梯度的表达式, 编码过程中非常容易出错, 训练结果难以移植到生产系统中……
\end{frame}

% \begin{frame}[fragile]{\insertsection}{\insertsubsection}
% \begin{columns}
% \column{\dimexpr\paperwidth-10pt}\vspace{-20pt}
% \begin{tcolorbox}[%
%     colback = barcolor!5!white,
%     colframe = headercolor,
%     left = 5mm,
%     top = 0pt,
%     bottom = 5pt,
%     enhanced,
% ]
% \begin{multicols}{1}
%     \begin{minted}[%
%         fontsize = \fontsize{8}{9}\selectfont,
%         linenos,
%         numbersep = 3mm,
%         autogobble,
%         breaklines,
%         breakafter = d
%     ]{python}
% from math import exp, log

% def sigmoid(x, n = 1):
%   return 1 / (1 + exp(-n * x)) if x >= 0 else exp(n * x) / (1 + exp(n * x))

% def l(x1, x2, y, w1, w2):
%   return -sum([y * log(sigmoid(z(x1, x2, w1, w2))) + (1 - y) * log(1 - sigmoid(z(x1, x2, w1, w2))) for x1, x2, y in zip(x1, x2, y)])

% def z(x1, x2, w1, w2):
%   return w1 * x1 + w2 * x2 + 1

% def get_delta_w1(x1, x2, y, w1, w2):
%   return -sum([(y - sigmoid(z(x1, x2, w1, w2))) * x1 for x1, x2, y in zip(x1, x2, y)])

% def get_delta_w2(x1, x2, y, w1, w2):
%   return -sum([(y - sigmoid(z(x1, x2, w1, w2))) * x2 for x1, x2, y in zip(x1, x2, y)])
%     \end{minted}
%   \end{multicols}
% \end{tcolorbox}
% \end{columns}
% \end{frame}

%!TEX root = ../main.tex

\section{TensorFlow 基本概念}
\subsection{张量 Tensor}
\begin{frame}{\insertsection}{\insertsubsection}
张量最初是个物理学的概念,
\begin{quote}
    A tensor is something that transforms like a tensor\\
    一个量, 在不同的参考系下按照某种特定的法则进行变换, 就是张量.\\
    \rule{0pt}{0pt}\hfill  --- A. Zee, "Einstein Gravity in a Nutshell"
\end{quote}

在 TensorFlow 中, 可以将张量理解成多维数组, 张量穿梭在数据流图中, 当经过节点时, 会被节点按照特定法则变换成另一个张量, 这个就是 TensorFlow 名字的由来.

张量的维数被称为张量的阶, 我们常见的标量, 其实是 $0$ 阶张量, 矢量就是 $1$ 阶张量, 矩阵就是 $2$ 阶张量, 一幅图片, 在 TensorFlow 中可以是 $3$ 阶张量.
\end{frame}

\begin{frame}{\insertsection}{\insertsubsection}
TensorFlow 文档中使用了三种记号来方便地描述张量的维度: 阶, 形状以及维数. 下表展示了他们之间的关系:\vspace{10pt}

\begin{table}
  \centering
  \begin{tabu}{clcl}
  \tabucline[1pt]{-}
  \rowfont{\bfseries}
  阶       & 形状                                     & 维数     & python 代码\\
  \hline
  $0$      & \texttt{[ ]}                             & $0$      & \inlinepython{t0 = 483}\\
  $1$      & \texttt{[$D_0$]}                         & $1$      & \inlinepython{t1 = [1.1, 2.2, 3.3]}\\
  $2$      & \texttt{[$D_0$, $D_1$]}                  & $2$      & \inlinepython{t2 = [[1, 2, 3], [4, 5, 6], [7, 8, 9]]}\\
  $3$      & \texttt{[$D_0$, $D_1$, $D_2$]}           & $3$      & \inlinepython{t3 = [[[1], [2]], [[3], [4]], [[5], [6]]]}\\
  $\vdots$ & $\vdots$                                 & $\vdots$ & $\vdots$\\
  $n$      & \texttt{[$D_0$, $D_1$, $\cdots$, $D_n$]} & $n$      & \inlinepython{tn = ...}\\
  \tabucline[1pt]{-}
  \end{tabu}
\end{table}
\end{frame}

\begin{frame}{\insertsection}{\insertsubsection}
除了维度, Tensors有一个数据类型属性. 可为一个张量指定下列数据类型的任意一个类型:\vspace{10pt}

\begin{table}
  \centering
  \begin{tabu}{ll|ll}
  \tabucline[1pt]{-}
  \rowfont{\bfseries}
  类型 & 描述 & 类型 & 描述 \\
  \hline
  \inlinepython{tf.float32} & 32 位浮点数     & \inlinepython{tf.string}    & 可变长度的字节数组\\
  \inlinepython{tf.float64} & 64 位浮点数     & \inlinepython{tf.bool}      & 布尔型\\
  \inlinepython{tf.int64}   & 64 位有符号整型 & \inlinepython{tf.complex64} & 由两个 32 位浮点数组成的复数\\
  \inlinepython{tf.int32}   & 32 位有符号整型 & \inlinepython{tf.qint32}    & 用于量化操作的 32 位有符号整型\\
  \inlinepython{tf.int16}   & 16 位有符号整型 & \inlinepython{tf.qint8}     & 用于量化操作的 8 位有符号整型\\
  \inlinepython{tf.int8}    & 8 位有符号整型  & \inlinepython{tf.quint8}    & 用于量化操作的 8 位无符号整型\\
  \inlinepython{tf.uint8}   & 8 位无符号整型  & & \\
  \tabucline[1pt]{-}
  \end{tabu}
\end{table}
\end{frame}

\begin{frame}[fragile]{\insertsection}{\insertsubsection}
TensorFlow 中张量包括 \inlinepython{constant}, \inlinepython{placeholder} 和 \inlinepython{Variable}.

\begin{onlyenv}<2>
\inlinepython{tf.constant()} 函数提供在 TensorFlow 中定义不可更改张量的方法, 定义如下所示.

\begin{pythoncode}
def constant(value, dtype=None, shape=None, name="Const", verify_shape=False)
\end{pythoncode}

\begin{itemize}
\item \inlinepython{value}, 符合 TensorFlow 中定义的数据类型的常数值或者常数列表;
\item \inlinepython{dtype}, 数据类型, 可选;
\item \inlinepython{shape}, 常量的形状, 可选;
\item \inlinepython{name}, 常量的名字, 可选;
\item \inlinepython{verify_shape}, 常量的形状是否可以被更改, 默认不可更改.
\end{itemize}

除了直接赋值以外, 还可使用 \inlinepython{tf.ones()}, \inlinepython{tf.zeros()} 等初始化张量的方法.
\end{onlyenv}

\begin{onlyenv}<3>
\inlinepython{tf.placeholder()} 函数提供在 TensorFlow 中定义占位张量的方法, 定义如下所示.

\begin{pythoncode}
def placeholder(dtype, shape=None, name=None)
\end{pythoncode}

\begin{itemize}
\item \inlinepython{dtype} 指定占位张量的数据类型, 可以是 TensorFlow 中的数据类型, 如常用的 \inlinepython{tf.float32}, \inlinepython{tf.float64} 等数值类型;
\item \inlinepython{shape} 表示数据类型, 默认的 \inlinepython{None} 是一个一维的数值, \inlinepython{shape = [None, 5]}, 表示行不定, 列是 \inlinepython{5};
\item \inlinepython{name} 是张量名称;
\end{itemize}

占位变量是一种 TensorFlow 用来解决读取大量训练数据问题的机制, 它允许你现在不用给它赋值, 随着训练的开始, 再把训练数据传送给训练网络学习.
\end{onlyenv}

\begin{onlyenv}<4->
TensorFlow 中变量是通过 \inlinepython{Variable} 类来实现的, 初始化函数 \inlinepython{tf.Variable()} 定义如下.

\begin{pythoncode}
def __init__(self, initial_value=None, trainable=True, collections=None, validate_shape=True, caching_device=None, name=None, variable_def=None, dtype=None, expected_shape=None, import_scope=None)
\end{pythoncode}


\only<4>{
\begin{itemize}
\item \inlinepython{initial_value}, 初始值, 必填, 张量或可以转换为张量的 Python 对象;
\item \inlinepython{trainable}, 如果参数 \inlinepython{trainable} 的值为 \inlinepython{True}, 则默认值也将变量添加到图形中集合 \inlinepython{GraphKeys.TRAINABLE_VARIABLES}。这个集合用作 \inlinepython{Optimizer} 类使用的默认变量列表;
\item \inlinepython{collections}, 新的变量被添加到这些集合。默认为 \inlinepython{[GraphKeys.GLOBAL_VARIABLES]};
\end{itemize}
}
\only<5>{
\begin{itemize}
\item \inlinepython{validate_shape}, 如果 \inlinepython{False}, 允许变量用初始化未知形状的值。如果 \inlinepython{True}, 默认的形状 \inlinepython{initial_value} 必须是已知的;
\item \inlinepython{caching_device}, 指定变量的默认缓存位置;
\item \inlinepython{name}, 是张量名称;
\end{itemize}
}
\only<6>{
\begin{itemize}
\item \inlinepython{variable_def}, 指定 \inlinepython{VariableDef} 的协议缓冲区;
\item \inlinepython{dtype}, 指定张量的数据类型;
\item \inlinepython{expected_shape}, 是张量的形状, 如果设置, \inlinepython{initial_value} 需要符合这个形状;
\item \inlinepython{import_scope}, 张量名称前缀.
\end{itemize}
}

\end{onlyenv}
\end{frame}

\subsection{节点 Operator}
\begin{frame}{\insertsection}{\insertsubsection}
在 TensorFlow 中, 每一个节点代表着一个操作, 一般用来表示施加的的数学运算, 也可以表示数据输入的起点以及输出的终点. 下面是一些重要的操作:

\begin{table}
  \centering
  \begin{tabu}{ll|ll}
  \tabucline[1pt]{-}
  \rowfont{\bfseries}
   操作 & 描述 &  操作 & 描述 \\
  \hline
  \inlinepython{tf.add(x, y, name=None)} & 求和 & \inlinepython{tf.sign(x, name=None)}   & 返回符号\\
  \inlinepython{tf.sub(x, y, name=None)} & 减法 & \inlinepython{tf.neg(x, name=None)}    & 取负 ($y = -x$)\\
  \inlinepython{tf.mul(x, y, name=None)} & 乘法 & \inlinepython{tf.square(x, name=None)} & 计算平方 ($y = x^2$)\\
  \inlinepython{tf.div(x, y, name=None)} & 除法 & \inlinepython{tf.round(x, name=None)}  & 求最接近的整数\\
  \inlinepython{tf.mod(x, y, name=None)} & 取模 & \inlinepython{tf.sqrt(x, name=None)}   & 开根号 ($y = \sqrt{x}$)\\
  \inlinepython{tf.pow(x, y, name=None)} & 乘幂 & \inlinepython{tf.abs(x, name=None)}    & 绝对值\\
  \inlinepython{tf.inv(x, name=None)}    & 取反 & \inlinepython{tf.exp(x, name=None)}    & 计算 $\e$ 的次方\\
  \tabucline[1pt]{-}
  \end{tabu}
\end{table}
\end{frame}

\subsection{数据流图 Graph}
\begin{frame}{\insertsection}{\insertsubsection}
\vspace{10pt}
\begin{minipage}[m]{0.5\textwidth}
    TensorFlow 是用数据流图对计算过程进行描述的. 在数据流图中, 节点代表数学运算, 边表示节点之间的某种联系, 负责在节点之间传输即张量.

    节点可以被分配到多个计算设备上, 可以异步和并行的进行操作. 因为是有向图, 所以只能等待之前的节点运行结束, 当前节点才能执行操作.

    如图所示是一个简单的计算流图.%
    \end{minipage}\hfill
    \begin{minipage}[m]{0.4\textwidth}
    \inlineframe{\scalebox{0.8}{%!TEX root = ../main.tex

\begin{tikzpicture}[%
    node/.style = {circle, draw, minimum size = 0.8cm, align = flush center, inner sep = 0pt, fill = black, text = white},
    y = 0.8cm]

    \node[node] (x) at (5, 0) {$x$};
    \node[node] (11) at (6, 0) {$11$};
    \node[node] (y) at (9, 0) {$y$};
    \node[node] (div) at (5.5, 2) {$\div$};
    \node[node] (22) at (3.2, 2) {$2.2$};
    \node[node] (7) at (7.5, 2) {$7$};
    \node[node] (cos) at (9.5, 2) {$\cos$};
    \node[node] (sub) at (4.5, 4) {$-$};
    \node[node] (times) at (8.5, 4) {$\times$};
    \node[node] (plus) at (6.5, 6) {$+$};

    \draw[->] (x) -- (div);
    \draw[->] (11) -- (div);
    \draw[->] (y) -- (cos);
    \draw[->] (22) -- (sub);
    \draw[->] (div) -- (sub);
    \draw[->] (7) -- (times);
    \draw[->] (cos) -- (times);
    \draw[->] (sub) -- (plus);
    \draw[->] (times) -- (plus);
\end{tikzpicture}
}}
    \[
        \text{表达式: }\bigg(2.2 - \Big(\frac{x}{11}\Big)\bigg) + \big(7 \times \cos(y)\big)\text{.}
    \]
\end{minipage}
\end{frame}

\begin{frame}[fragile]{\insertsection}{\insertsubsection}
\begin{pythoncode}
import tensorflow as tf
import numpy as np

x1 = tf.placeholder(tf.float32, [None, 1])
w1 = tf.Variable(tf.random_normal([1, 10]))
b1 = tf.Variable(tf.ones([1, 10]))

x2 = tf.nn.relu(tf.matmul(x1, w1) + b1)
w2 = tf.Variable(tf.random_normal([10, 1]))
b2 = tf.Variable(tf.ones([1, 1]))

predict_values = tf.matmul(x2, w2) + b2
actual_values = tf.placeholder(tf.float32, [None, 1])
loss = tf.reduce_mean(tf.reduce_sum(tf.square(actual_values - predict_values), reduction_indices = [1]))
\end{pythoncode}

\vspace{-7.5cm}\hspace{10.5cm}\scalebox{0.53}{\begin{tikzpicture}[%
  tensor/.style = {rectangle, inner sep = 2pt, rounded corners = 0.4cm, draw, align = flush center, minimum height = 0.8cm, minimum width = 0.8cm},
  operator/.style = {tensor, fill = black, draw = none, text = white, inner sep = 2pt},
  line width = 1pt,
  y = 1.1cm
]
\node[tensor, circle]   (x1) at (9, 0) {$x_1$};
\node[tensor, circle]   (w1) at (7, 0) {$W_1$};
\node[operator]         (t1) at (8, 1) {\mintinline{text}{tf.matmul}};
\node[tensor, circle]   (b1) at (6, 1) {$b_1$};
\node[operator]         (p1) at (7, 2) {\mintinline{text}{tf.add}};
\node[operator]         (r1) at (7, 3) {\mintinline{text}{tf.nn.relu}};
\node[tensor, circle]   (w2) at (5, 3) {$W_2$};
\node[operator]         (t2) at (6, 4) {\mintinline{text}{tf.matmul}};
\node[tensor, circle]   (b2) at (4, 4) {$b_2$};
\node[operator]         (p2) at (5, 5) {\mintinline{text}{tf.add}};
\node[tensor]           (pv) at (5, 6) {\mintinline{text}{prediction_values}};
\node[tensor]           (av) at (8, 6) {\mintinline{text}{actual_values}};
\node[operator]         (sb) at (6.5, 7.5) {\mintinline{text}{tf.sub}};
\node[operator]         (sq) at (6.5, 8.5) {\mintinline{text}{tf.square}};
\node[operator]         (rs) at (6.5, 9.5) {\mintinline{text}{tf.reduce_sum}};
\node[operator]         (rm) at (6.5, 10.5) {\mintinline{text}{tf.reduce_mean}};
\node[tensor]           (ls) at (6.5, 11.5) {\mintinline{text}{loss}};

\draw[->] (x1) -- (t1);
\draw[->] (w1) -- (t1);
\draw[->] (b1) -- (p1);
\draw[->] (t1) -- (p1);
\draw[->] (p1) -- (r1);
\draw[->] (w2) -- (t2);
\draw[->] (r1) -- (t2);
\draw[->] (b2) -- (p2);
\draw[->] (t2) -- (p2);
\draw[->] (p2) -- (pv);
\draw[->] (pv) -- (sb);
\draw[->] (av) -- (sb);
\draw[->] (sb) -- (sq);
\draw[->] (sq) -- (rs);
\draw[->] (rs) -- (rm);
\draw[->] (rm) -- (ls);
\end{tikzpicture}}
\end{frame}

\subsection{会话 Session}
\begin{frame}[fragile]{\insertsection}{\insertsubsection}
构造阶段完成后, 才能启动图。启动图的第一步是创建一个 Session 对象, 如果无任何创建参数, 会话构造器将启动默认图。会话会管理TensorFlow程序运行时的所有资源。当所有计算完成之后需要关闭会话来帮助系统回收资源，否则就可能出现资源泄露的问题。

\begin{minipage}[t]{0.48\textwidth}
\begin{pythoncode}
import tensorflow as tf

a = tf.constant([1.0, 2.0])
b = tf.constant([2.0, 3.0])
result = a + b

session = tf.Session()
print(session.run(result))
session.close()
\end{pythoncode}
\end{minipage}%
\hfill%
\begin{minipage}[t]{0.48\textwidth}
\begin{pythoncode}
import tensorflow as tf

a = tf.constant([1.0, 2.0])
b = tf.constant([2.0, 3.0])
result = a + b

with tf.Session() as session:
    print(session.run(result))
# session.close()
\end{pythoncode}
\end{minipage}
\end{frame}
%!TEX root = ../main.tex

\section{TensorFlow 实战}
\begin{frame}{\secname}
\vfill\hspace{20pt}
\begin{minipage}[m]{0.25\textwidth}
  \linespread{1}\large
  \tableofcontents[sectionstyle = hide/hide, subsectionstyle = show/show/hide]
\end{minipage}%
\hfill%
\begin{minipage}[m]{0.6\textwidth}
  \includegraphics[height = 7cm]{./icons/programming.pdf}
\end{minipage}%
\vfill
\end{frame}

\subsection{分类问题}
\begin{frame}{\insertsection}{\insertsubsection}
\begin{minipage}[m]{0.48\textwidth}
回到最开始的那个二分类问题, 看看用 \tensorflow{} 是如何解决这个问题的, $500$ 个点的分布如下图所示.

我们通过观察数据的特征, 仍然选用一样的神经网络, 如下图所示.

\centering
\scalebox{0.8}{%
\begin{tikzpicture}[%
  node/.style = {circle, draw, inner sep = 0pt, minimum size = 0.8cm},
  x = 1.4cm]
  \node[node] (x1) at (0, 2) {$x_1$};
  \node[node] (x2) at (0, 1) {$x_2$};
  \node[node] (b)  at (0, 0) {$1$};
  \node[node] (sum) at (2, 1) {$\sum$};
  \node[node] (step) at (4, 1) {$H(x)$};

  \draw[->] (x1) -- (sum) node[midway, fill = white, inner sep = 1pt, sloped] {$W_1$};
  \draw[->] (x2) -- (sum) node[midway, fill = white, inner sep = 1pt, sloped] {$W_2$};
  \draw[->] (b) -- (sum);
  \draw[->] (sum) -- (step);
  \draw[->] (step) -- ++(0:1.5cm) node[anchor = west] {$y$};
\end{tikzpicture}}
\end{minipage}%
\hfill%
\begin{minipage}[m]{0.48\textwidth}
\begin{tikzpicture}
  \begin{axis}[%
    xlabel = $x_1$,
    ylabel = $x_2$,
    height = 7.8cm,
    width = 7cm,
    title = $500$ 个点的分布图
  ]
    \addplot[only marks, discard if not={type}{1}, blue] table[col sep=comma, x = x1, y = x2]{./data/classification/training_data.dat};
    \addplot[only marks, discard if not={type}{0}, green] table[col sep=comma, x = x1, y = x2]{./data/classification/training_data.dat};
  \end{axis}
\end{tikzpicture}
\end{minipage}
\end{frame}

\begin{frame}[fragile]{\insertsection}{\insertsubsection}
这次我们不用计算神经网络的表达式, 直接照着神经网络写 \tensorflow{} 代码即可.

\begin{pythoncode}[fontsize = \fontsize{8}{8}\selectfont]
import tensorflow as tf
import numpy as np

# 定义 $\bm{x} = (x_1, x_2)$, $\bm{W} = (W_1, W_2)$ 和 $b = 1$
x = tf.placeholder(tf.float32, [None, 2])
w = tf.Variable(tf.zeros([2, 1]))
b = tf.constant([1.])

# 定义神经网络输出
predict_type = tf.sigmoid((tf.matmul(x, w) + b) * 10)
# 定义真实输出 $y = g(\bm{W}\bm{x} + 1)$
actual_type = tf.placeholder(tf.float32, [None, 1])
# 定义目标函数 $\ell(\bm{W}) =  -\sum_{i = 1}^{500} y_i\ln g(\bm{W}\bm{x}_i + 1) + (1 - y_i)\ln\big(1 - g(\bm{W}\bm{x}_i + 1)\big)$
loss = -tf.reduce_sum(actual_type * tf.log(predict_type) + (1 - actual_type) * tf.log(1 - predict_type))
\end{pythoncode}

\vspace{-6.3cm}\hfill%
\scalebox{0.8}{%
\begin{tikzpicture}[%
  node/.style = {circle, draw, inner sep = 0pt, minimum size = 0.8cm},
  x = 1.4cm, y = 1.4cm]
  \node[node] (x1) at (0, 2) {$x_1$};
  \node[node] (x2) at (0, 1) {$x_2$};
  \node[node] (b)  at (0, 0) {$1$};
  \node[node] (sum) at (2, 1) {$\sum$};
  \node[node] (step) at (4, 1) {$H(x)$};

  \draw[->] (x1) -- (sum) node[midway, fill = barcolor!5!white, inner sep = 1pt, sloped] {$W_1$};
  \draw[->] (x2) -- (sum) node[midway, fill = barcolor!5!white, inner sep = 1pt, sloped] {$W_2$};
  \draw[->] (b) -- (sum);
  \draw[->] (sum) -- (step);
  \draw[->] (step) -- ++(0:1.5cm) node[anchor = west] {$y$};
\end{tikzpicture}\ \ }
\end{frame}

\begin{frame}[fragile]{\insertsection}{\insertsubsection}
\begin{pythoncode}[fontsize = \fontsize{8}{8}\selectfont]
def load_training_data(path): # 加载训练数据 $\bm{x}$ 和 $\bm{y}$
  with open(path, 'r', encoding = 'utf8') as file:
    next(file)
    lines = [[float(item) for item in line.split()] for line in file]
    xs = np.array([line[:2] for line in lines])
    types = np.array([line[2] for line in lines])[:, None]
  return xs, types

with tf.Session() as session:
  session.run(tf.global_variables_initializer()) # 初始化变量
  train_step = tf.train.GradientDescentOptimizer(0.003).minimize(loss) # 梯度下降法
  xs, types = load_training_data('./training.dat') # 加载训练数据 $\bm{x}$ 和 $\bm{y}$
  for i in range(1000): # 训练 $1000$ 次
    session.run(train_step, feed_dict = {x:xs, actual_type:types}) # 填充训练数据
    # 打印目标函数值
    if i % 10 == 0: print(session.run(loss, feed_dict = {x:xs, actual_type:types}))
  print(session.run(w, feed_dict = feed_dict)) # 打印最终结果
\end{pythoncode}
\end{frame}

\begin{frame}[fragile]{\insertsection}{\insertsubsection}
最终, 参数 $W_1$ 和 $W_2$ 收敛于点 $(1.07, -3.10)$, 也是非常理想的结果.

我们还可以换一个目标函数, 用模型预测错误的数量作为目标函数, 其表达式如下所示:
\begin{flalign*}
&& \ell(\bm{W}) &= \sum_{\mathclap{x_1, x_2, y}} \big|H(W_1x_1+ W_2x_2 + 1)) - y\big|\text{,} &&\\
\text{由于绝对值不可微,}&& \text{上式} &= \sum_{\mathclap{x_1, x_2, y}} \big(H(W_1x_1+ W_2x_2 + 1)) - y\big)^2\text{,} &&\text{\hphantom{由于绝对值不可微,}}\\
\text{Sigmoid 函数替换,}&& &\approx\sum_{\mathclap{x_1, x_2, y}} \big(g(W_1x_1+ W_2x_2 + 1)) - y\big)^2\text{.} &&
\end{flalign*}

我们修改目标函数的表达式, 参数 $W_1$ 和 $W_2$ 收敛于点 $(1.07, -3.07)$, 效果拔群.
\begin{pythoncode}
loss = tf.reduce_sum(tf.square(actual_type - predict_type))
\end{pythoncode}
\end{frame}

\subsection{拟合问题}
\begin{frame}{\insertsection}{\insertsubsection}
\begin{minipage}[t]{0.41\textwidth}
有 300 个点, 其分布如下图所示, 现在的目标是用一个神经网络, 来拟合这 300 个点.

\begin{center}
\hspace{-10pt}%
\begin{tikzpicture}
\begin{axis}[width = 1.2\textwidth,
             height = 6.4cm]
\addplot[only marks] table[col sep = comma] {./data/fitting/training_data.dat};
\end{axis}
\end{tikzpicture}
\end{center}
\end{minipage}%
\hfill%
\begin{minipage}[t]{0.55\textwidth}
我门用一个两层网络来实现对这些离散点的拟合, 这两层网络的公式为:%
%
\begin{align*}
\bm{x}_2 &= \relu(\bm{W}_1x_1 + \bm{b}_1)\text{,} && \text{$\bm{W}_1\in\mathbb{R}^{10\times1}$, $\bm{b}_1\in\mathbb{R}^{10\times1}$,}\\
\hat{y} &= \bm{W}_2\bm{x}_2 + b_2\text{,}  && \text{$\bm{W}_2\in\mathbb{R}^{1\times10}$, $b_2\in\mathbb{R}^{1\times1}$.}
\end{align*}%
其中, 函数 $\relu(\,\cdot\,)$ 被称为线性整流函数(Rectified Linear Unit, ReLU), 图像如下所示.\vspace{-2pt}

\begin{center}
\hspace{-10pt}%
\begin{tikzpicture}
  \begin{axis}[width = 1.2\textwidth,
               height = 3.7cm,
               legend pos = north west,
  ]
    \addplot[mark = none, domain = -5:2] {(x < 0) * 0 + (x >= 0) * x};
    \legend{{$\mathrm{relu}(x) = \max{(0, x)}$}};

  \end{axis}
\end{tikzpicture}
\end{center}
\end{minipage}
\end{frame}

\begin{frame}[fragile]{\insertsection}{\insertsubsection}
\begin{pythoncode}[fontsize = \fontsize{8}{8}\selectfont, escapeinside=||]
import tensorflow as tf
import numpy as np

x1 = tf.placeholder(tf.float32, [None, 1])
w1 = tf.Variable(tf.random_normal([1, 10]))
b1 = tf.Variable(tf.ones([1, 10]))
# $\bm{x}_2 = \relu(\bm{W}_1x_1 + \bm{b}_1)$, $\bm{W}_1\in\mathbb{R}^{10\times1}$, $\bm{b}_1\in\mathbb{R}^{10\times1}$
x2 = tf.nn.relu(tf.matmul(x1, w1) + b1)
w2 = tf.Variable(tf.random_normal([10, 1]))
b2 = tf.Variable(tf.ones([1, 1]))
# $\hat{y} = \bm{W}_2\bm{x}_2 + b_2$, $\bm{W}_2\in\mathbb{R}^{1\times10}$, $b_2\in\mathbb{R}^{1\times1}$
predict_values = tf.matmul(x2, w2) + b2 # 预测值
actual_values = tf.placeholder(tf.float32, [None, 1]) # 真实值

# 目标函数 $\ell(\bm{W}_1, \bm{b}_1, \bm{W}_2, b_2) = \frac{1}{n}\sum_{i = 1}^{300}(y - \hat{y})^2$
loss = tf.reduce_mean(tf.reduce_sum(tf.square(actual_values - |\hspace{50pt}| predict_values), reduction_indices = [1]))
\end{pythoncode}

\vspace{-7.5cm}\hspace{10.5cm}\scalebox{0.53}{\begin{tikzpicture}[%
  tensor/.style = {rectangle, inner sep = 2pt, rounded corners = 0.4cm, draw, align = flush center, minimum height = 0.8cm, minimum width = 0.8cm},
  operator/.style = {tensor, fill = black, draw = none, text = white, inner sep = 2pt},
  line width = 1pt,
  y = 1.1cm
]
\node[tensor, circle]   (x1) at (9, 0) {$\bm{x}_1$};
\node[tensor, circle]   (w1) at (7, 0) {$\bm{W}_1$};
\node[operator]         (t1) at (8, 1) {\mintinline{text}{tf.matmul}};
\node[tensor, circle]   (b1) at (6, 1) {$\bm{b}_1$};
\node[operator]         (p1) at (7, 2) {\mintinline{text}{tf.add}};
\node[operator]         (r1) at (7, 3) {\mintinline{text}{tf.nn.relu}};
\node[tensor, circle]   (w2) at (5, 3) {$\bm{W}_2$};
\node[operator]         (t2) at (6, 4) {\mintinline{text}{tf.matmul}};
\node[tensor, circle]   (b2) at (4, 4) {$b_2$};
\node[operator]         (p2) at (5, 5) {\mintinline{text}{tf.add}};
\node[tensor]           (pv) at (5, 6) {\mintinline{text}{prediction_values}};
\node[tensor]           (av) at (8, 6) {\mintinline{text}{actual_values}};
\node[operator]         (sb) at (6.5, 7.5) {\mintinline{text}{tf.sub}};
\node[operator]         (sq) at (6.5, 8.5) {\mintinline{text}{tf.square}};
\node[operator]         (rs) at (6.5, 9.5) {\mintinline{text}{tf.reduce_sum}};
\node[operator]         (rm) at (6.5, 10.5) {\mintinline{text}{tf.reduce_mean}};
\node[tensor]           (ls) at (6.5, 11.5) {\mintinline{text}{loss}};

\draw[->] (x1) -- (t1);
\draw[->] (w1) -- (t1);
\draw[->] (b1) -- (p1);
\draw[->] (t1) -- (p1);
\draw[->] (p1) -- (r1);
\draw[->] (w2) -- (t2);
\draw[->] (r1) -- (t2);
\draw[->] (b2) -- (p2);
\draw[->] (t2) -- (p2);
\draw[->] (p2) -- (pv);
\draw[->] (pv) -- (sb);
\draw[->] (av) -- (sb);
\draw[->] (sb) -- (sq);
\draw[->] (sq) -- (rs);
\draw[->] (rs) -- (rm);
\draw[->] (rm) -- (ls);
\end{tikzpicture}}
\end{frame}

\begin{frame}[fragile]{\insertsection}{\insertsubsection}
\begin{pythoncode}[fontsize = \fontsize{8}{8}\selectfont]
def load_training_data(path): # 读取训练数据
    with open(path, 'r', encoding = 'utf8') as file:
        next(file)
        lines = [[float(item) for item in line.split(',')] for line in file]
    return (np.array([line[i] for line in lines])[:, None] for i in [0, 1])

with tf.Session() as session: # 初始化 Session
  session.run(tf.global_variables_initializer()) # 初始化变量

  (input_data, output_data) = load_training_data('./training_data.dat') # 加载训练数据
  train_step = tf.train.GradientDescentOptimizer(0.1).minimize(loss) # 设置目标函数

  for i in range(1000): # 训练 $1000$ 次
    session.run(train_step, feed_dict = {x1:input_data, actual_values:output_data})
    if i % 50 == 0:
      print(session.run(loss, feed_dict = {x1:input_data, actual_values:output_data}))
\end{pythoncode}
\end{frame}

\begin{frame}{\insertsection}{\insertsubsection}
训练 $1000$ 次后, 目标函数 $\ell(\bm{W}_1, \bm{b}_1, \bm{W}_2, b_2) = 2.48\times 10^{-3}$, 此时拟合函数的图像如图所示.\vspace{10pt}

\begin{tikzpicture}
\begin{axis}[width = 1\textwidth,
             height = 6.4cm,
             xlabel = $x$,
             ylabel = $y$,
             title = 训练数据与拟合曲线,
             legend pos = south east]
\addplot[only marks] table[col sep = comma] {./data/fitting/training_data.dat};
\addplot[mark = none,white, line width = 1pt, double = red, double distance = 2pt, on layer = foreground] table[col sep = comma, x = x, y = y] {./data/fitting/fitting_line.dat};
\legend{训练数据, 拟合曲线}
\end{axis}
\end{tikzpicture}
\end{frame}

% \begin{frame}[fragile]{\insertsection}{\insertsubsection}
% 我门用一个两层网络来实现对这些离散点的拟合, 这两层网络的公式为:%
% %
% \begin{flalign*}
%   && x_2 &= \relu(W_1x_1 + b_1)\text{,} && \text{$W_1\in\mathbb{R}^{10\times1}$, $b_1\in\mathbb{R}^{10\times1}$,} &&\\
%   && \hat{y} &= W_2x_2 + b_2\text{,}                  && \text{$W_2\in\mathbb{R}^{1\times10}$, $b_2\in\mathbb{R}^{1\times1}$.} &&
% \end{flalign*}

% \begin{minipage}[m]{0.38\textwidth}
% \begin{tikzpicture}
%   \begin{axis}[width = 1.2\textwidth,
%                height = 5.1cm,
%                legend pos = north west,
%   ]
%     \addplot[mark = none, domain = -2:5] {(x < 0) * 0 + (x >= 0) * x};
%     \legend{{$\mathrm{relu}(x) = \max{(0, x)}$}};

%   \end{axis}
% \end{tikzpicture}
% \end{minipage}%
% \hfill%
% \begin{minipage}[m]{0.60\textwidth}
% \begin{pythoncode}
% x1 = tf.placeholder(tf.float32, [None, 1])
% w1 = tf.Variable(tf.random_normal([1, 10]))
% b1 = tf.Variable(tf.ones([1, 10]))

% x2 = tf.nn.relu(tf.matmul(x1, w1) + b1)
% w2 = tf.Variable(tf.random_normal([10, 1]))
% b2 = tf.Variable(tf.ones([1, 1]))
% y = tf.matmul(x2, w2) + b2
% \end{pythoncode}
% \end{minipage}
% \end{frame}

% \subsection{训练网络}
% \begin{frame}[fragile]{\insertsection}{\insertsubsection}
% 我们用统计学中的均方误差作为代价函数来训练网络, 如下所示.%
% %
% \[
%   \ell(W_1, b_1, W_2, b_2) = \frac{1}{n}\sum_{i = 1}^n (y - \hat{y})^2\text{.}
% \]%
% %
% 根据上式可以写出 TensorFlow 中的训练目标.
% \begin{pythoncode}
% actual_values = tf.placeholder(tf.float32, [None, 1])
% loss = tf.reduce_mean(tf.reduce_sum(tf.square(actual_values - y), reduction_indices = [1]))

% train_step = tf.train.GradientDescentOptimizer(0.1).minimize(loss)
% \end{pythoncode}

% 我们可以看出, \inlinepython{tf.placeholder} 一般对应着训练数据, 因为训练数据只能在训练的时候给出, 在定义计算图的时候不知道具体的值, \inlinepython{tf.placeholder} 起占位的作用.
% \end{frame}

% \begin{frame}[fragile]{\insertsection}{\insertsubsection}
% 训练网络的过程, 就是求使代价函数 $\ell(W_1, b_1, W_2, b_2)$ 达到最小值时的参数集合 $W_1$, $b_1$, $W_2$, $b_2$. 而且需要启动一个会话 \inlinepython{tf.Session} 来完成运算图的计算.

% \begin{pythoncode}
% session = tf.Session()
% session.run(tf.initialize_all_variables())

% for step in range(1000):
%     feed_dict = {x1:input_data, actual_values:output_data}
%     session.run(train_step, feed_dict = feed_dict)

%     if step % 50 == 0:
%         print(session.run(loss, feed_dict = feed_dict))

% session.close()
% \end{pythoncode}
% \end{frame}

% \subsection{测试网络}

% \subsection{}
%!TEX root = ../main.tex

\section{总结}
\begin{frame}{\insertsection}
\vfill
\begin{minipage}[m]{0.55\textwidth}
\begin{itemize}
\item 以最直观的方式帮助大家理解神经网络, 目标函数, 梯度下降等概念;
\item 介绍了机器学习的一般步骤:
\begin{itemize}
\item 数据分析以及预处理,
\item 选择合适的模型,
\item 设定目标函数以及训练参数,
\item 验证以及调整模型;
\end{itemize}
\item 介绍了 \tensorflow{} 中张量, 节点, 数据流图, 会话等基本概念概念;
\item 介绍了 \tensorflow{} 如何解决机器学习中分类和拟合这两个经典问题.
\end{itemize}
\end{minipage}%
\hfill%
\begin{minipage}[m]{0.41\textwidth}
\includegraphics[width = \textwidth]{icons/conclusion.pdf}
\end{minipage}
\vfill
\end{frame}

\insertthanks

\end{document}
